\chapter{Generalità dell'anatomia}

% ---------------------------------------------------------------------------
\section{L'apparato scheletrico}

L'\textbf{apparato scheletrico} svolge funzioni di protezione, sostegno, contributo al movimento, riserva di minerali ed emopoiesi. Si distingue in:

\begin{itemize}
  \item \textbf{Scheletro assile}: cranio, vertebre, coste, sterno, sacro, cartilagini e legamenti. Protegge l'encefalo, il midollo spinale, gli organi di senso e i tessuti molli della cavità toracica; sostiene il peso corporeo sopra gli arti inferiori.
  \item \textbf{Scheletro appendicolare}: arti, ossa di sostegno e legamenti. Fornisce sostegno interno e posizione agli arti; sostiene e sposta lo scheletro assile.
\end{itemize}

Lo scheletro costituisce circa il \textbf{20\%} del peso corporeo; nell'adulto è composto da circa \textbf{200 ossa}. Il \textbf{midollo osseo} è il sito primario di produzione dei globuli rossi e bianchi del sangue.

% ---------------------------------------------------------------------------
\section{Struttura dell'osso}

\subsection{Composizione e tipi di tessuto osseo}

L'osso è costituito da una \textbf{matrice extracellulare} molto dura, compatta e mineralizzata, e da una componente cellulare con abbondanza di \textbf{osteociti}.

Si distinguono due tipi principali di tessuto osseo:

\begin{itemize}
  \item \textbf{Tessuto osseo spugnoso}: forma una rete di lamine e trabecole; riempie le estremità delle ossa lunghe e forma lo strato medio delle ossa piatte (es.\ lo sterno).
  \item \textbf{Tessuto osseo compatto}: relativamente denso e solido; tessuto densamente calcificato che forma la superficie esterna di tutte le ossa.
\end{itemize}

\subsection{Classificazione in base alla forma}

Le ossa si classificano in cinque categorie morfologiche:

\begin{itemize}
  \item \textbf{Ossa lunghe}: si sviluppano in lunghezza (es.\ radio, ulna, omero, tibia, perone).
  \item \textbf{Ossa corte}: forma cuboidale (es.\ ossa del carpo e del tarso).
  \item \textbf{Ossa piatte}: lunghezza e larghezza superiori allo spessore (es.\ scapole, sterno, ossa del cranio).
  \item \textbf{Ossa irregolari}: non classificabili nelle categorie precedenti (es.\ vertebre).
  \item \textbf{Ossa sesamoidi}: incluse all'interno di tendini (es.\ patella, osso pisiforme).
\end{itemize}

\subsection{Il midollo osseo}

Il \textbf{midollo osseo} è un tessuto molle che occupa i canali delle ossa lunghe e la fascia centrale delle ossa piatte. Contiene:

\begin{enumerate}
  \item \textbf{Midollo giallo}: costituito prevalentemente da stroma adiposo (adipociti).
  \item \textbf{Midollo rosso}: parenchima emopoietico contenente cellule del sangue rosse e bianche (mature e immature) e cellule staminali.
\end{enumerate}

\subsection{Vascolarizzazione}

L'osso è riccamente vascolarizzato e innervato. Le due fonti di irrorazione sono:

\begin{enumerate}
  \item Le \textbf{arterie nutritizie}
  \item Le \textbf{arterie periostee}
\end{enumerate}

La vascolarizzazione diafisaria e metaepifisaria è fittamente anastomizzata tra loro.

\subsection{Caratteristiche di superficie e periostio}

Le superfici esterne delle ossa sono rivestite da una spessa capsula connettivale costituita da tessuto connettivo denso a fibre intrecciate: il \textbf{periostio}. Il periostio protegge l'osso e supporta l'azione trofica mediata dai vasi sanguigni di cui è ricco.

Le principali caratteristiche di superficie delle ossa sono:

\resizebox{\linewidth}{!}{%
\begin{tabular}{@{}lll@{}}
\toprule
\textbf{Descrizione generale} & \textbf{Termine} & \textbf{Definizione ed esempio} \\
\midrule
Elevazioni e proiezioni (generale)
  & Processo & Proiezione o rilievo \\
  & Ramo & Parte di un osso che forma un angolo con il resto della struttura \\
\addlinespace
Processi nei punti di inserzione di tendini e legamenti
  & Trocantere & Voluminosa proiezione rugosa \\
  & Tuberosità & Proiezione rugosa più piccola \\
  & Tubercolo & Piccola proiezione rotondeggiante \\
  & Cresta & Linea rilevata \\
  & Linea & Linea sottile \\
  & Spina & Processo appuntito \\
\addlinespace
Processi per articolazione con ossa vicine
  & Testa & Estremità articolare espansa di un'epifisi, separata dalla diafisi da uno stretto collo \\
  & Collo & Più stretta connessione tra epifisi e diafisi \\
  & Condilo & Processo articolare liscio e tondeggiante \\
  & Troclea & Processo articolare liscio e scanalato, a forma di puleggia \\
  & Faccetta & Piccola superficie articolare piatta \\
\addlinespace
Depressioni
  & Fossa & Depressione poco profonda \\
  & Solco & Piccola fossa \\
\addlinespace
Aperture
  & Forame & Foro per il passaggio di vasi sanguigni e/o nervi \\
  & Fessura & Fenditura allungata \\
  & Meato o canale & Ingresso a un canale che attraversa a tutto spessore un osso \\
  & Seno o antro & Camera all'interno di un osso, in genere contenente aria \\
\bottomrule
\end{tabular}}

% ---------------------------------------------------------------------------
\section{Le articolazioni}

Le \textbf{articolazioni} rendono le ossa solidali e consentono il movimento reciproco delle ossa contigue, e quindi di segmenti scheletrici tra loro. In base alla mobilità si distinguono articolazioni \textbf{mobili} (es.\ spalla), \textbf{semimobili} (es.\ tra le vertebre) e \textbf{fisse} (es.\ ossa del cranio).

\subsection{Classificazione delle articolazioni}

\resizebox{\linewidth}{!}{%
\begin{tabular}{@{}llll@{}}
\toprule
\textbf{Cat.\ funzionale} & \textbf{Cat.\ strutturale} & \textbf{Descrizione} & \textbf{Esempio} \\
\midrule
\textbf{Sinartrosi} (nessun movimento)
  & Fibrose --- Suture & Connessioni fibrose tra segmenti scheletrici con ampie superfici di incastro & Tra le ossa del cranio \\
  & Fibrose --- Gonfosi & Connessioni fibrose con inserzioni nell'alveolo & Legamenti periodontali \\
  & Cartilaginee --- Sincondrosi & Interposizione di un disco cartilagineo & Disco epifisario \\
  & Ossee --- Sinostosi & Ossificazione di un altro tipo di sinartrosi & Parti del cranio; linee epifisarie \\
\addlinespace
\textbf{Amfiartrosi} (piccoli movimenti)
  & Fibrose --- Sindesmosi & Connessione legamentosa & Tra tibia e fibula \\
  & Cartilaginee --- Sinfisi & Connessione tramite disco fibrocartilagineo & Tra le ossa dell'anca; tra vertebre adiacenti \\
\addlinespace
\textbf{Diartrosi} (ampi movimenti)
  & Sinoviali --- Monoassiale & Movimento su un piano & Gomito, caviglia \\
  & Sinoviali --- Biassiale & Movimenti su due piani & Coste, polso \\
  & Sinoviali --- Triassiale & Movimenti su tre piani & Spalla, anca \\
\bottomrule
\end{tabular}}

\subsection{Le diartrosi sinoviali}

Le diartrosi rivestono le articolazioni mobili. La capsula articolare è composta da:

\begin{itemize}
  \item \textbf{Membrana fibrosa}: costituita da fibre collagene che si continuano fino al periostio; spesso rinforzata da legamenti capsulari.
  \item \textbf{Membrana sinoviale}: fissata ai margini della cartilagine articolare; secerne il \textbf{liquido sinoviale} (filtrato del plasma). Il liquido sinoviale vascolarizza, nutre e lubrifica la cavità articolare; la sua viscosità è data dall'acido ialuronico disciolto nel plasma.
\end{itemize}

\subsection{Tipi di articolazioni sinoviali}

\textbf{Artrodia} --- Superfici articolari pianeggianti che permettono leggeri movimenti di scivolamento (es.\ tra i processi articolari delle vertebre).

\textbf{Enartrosi} --- I due capi ossei sono ``sferici'', uno concavo e uno convesso. Permettono movimenti angolari su tutti i piani, inclusa la rotazione: tre assi di movimento, sei movimenti principali (es.\ articolazione coxo-femorale dell'anca, gleno-omerale della spalla).

\textbf{Condiloartrosi} --- Superfici articolari ellissoidali, una concava (cavità glenoidea) e una convessa (condilo). Permettono movimenti angolari su due piani: due assi di movimento, quattro movimenti principali (es.\ polso prossimale). L'\textbf{articolazione temporo-mandibolare} è una diartrosi doppia: due articolazioni sovrapposte con un disco completo interposto, una superiore (disco-fossa glenoide) e una inferiore (disco-condilo).

\textbf{Articolazione a sella} --- Le superfici hanno la forma di una sella di cavallo (concava longitudinalmente e convessa trasversalmente); permettono movimenti angolari su due piani e una rotazione assiale (es.\ tra il trapezio e il I osso metacarpale, articolazione femoro-rotulea).

\textbf{Ginglimo angolare (troclea)} --- La superficie articolare è a forma di puleggia, con l'asse perpendicolare alla diafisi; l'altra superficie è un'incavatura con una cresta corrispondente alla gola della troclea. Permette movimenti angolari su un piano: un asse, due movimenti principali (es.\ articolazione omero-ulnare, femoro-tibiale).

\textbf{Ginglimo laterale (trocoide)} --- Due capi ossei cilindrici, uno cavo e uno pieno, con l'asse parallelo all'asse longitudinale delle ossa. Permette un movimento rotatorio: un asse, due movimenti principali (es.\ articolazione prossimale radio-ulnare, articolazione atlanto-assiale mediana).

\subsection{I movimenti articolari}

I movimenti si classificano in: \textbf{lineare} (scivolamento, senza variazione angolare), \textbf{angolare} (variazione dell'angolo tra due segmenti), \textbf{circumduzione} (l'estremità distale descrive un cerchio) e \textbf{rotazione} (attorno all'asse longitudinale).

Termini specifici per gli arti: flessione/estensione, abduzione/adduzione, supinazione/pronazione, eversione/inversione, flessione dorsale/plantare, retrazione/protrusione, opposizione, depressione/elevazione.

% ---------------------------------------------------------------------------
\section{Orientamento e topografia corporea}

\subsection{Posizione anatomica}

La \textbf{posizione anatomica} è il riferimento convenzionale: stazione eretta, viso rivolto in avanti, arti paralleli all'asse del corpo, palmi delle mani rivolti in avanti (supinazione), piedi leggermente divaricati. I termini di posizione e i piani corporei si definiscono sempre in relazione a questa posizione.

\subsection{Termini di posizione}

I termini di posizione caratterizzano la situazione statica del corpo o di sue parti:

\begin{itemize}
  \item \textbf{Craniale} (o cefalico): vicino al polo cefalico.
  \item \textbf{Caudale} (o podalico): vicino al polo caudale.
  \item \textbf{Mediale}: vicino al piano di simmetria.
  \item \textbf{Laterale}: lontano dal piano di simmetria.
  \item \textbf{Ventrale}: posto anteriormente.
  \item \textbf{Dorsale}: posto posteriormente.
  \item \textbf{Prossimale}: vicino all'inserzione (negli arti).
  \item \textbf{Distale}: vicino all'estremità libera (negli arti).
\end{itemize}

\subsection{Piani e assi di simmetria}

I \textbf{piani} utilizzati per suddividere il corpo sono:

\begin{itemize}
  \item \textbf{Piani sagittali}: il piano sagittale che divide il corpo in due parti specularmente uguali (\textit{antimeri}) è detto \textbf{piano sagittale mediano}; i piani paralleli a esso si chiamano parasagittali.
  \item \textbf{Piani frontali}: andamento cranio-caudale, perpendicolari ai piani sagittali e trasversali; separano la porzione ventrale da quella dorsale. Non hanno un piano di riferimento fisso.
  \item \textbf{Piani trasversali}: andamento latero-laterale, paralleli alla superficie di appoggio in posizione eretta; dividono la parte craniale da quella caudale.
\end{itemize}

I tre assi principali sono: \textbf{craniocaudale}, \textbf{sagittale} e \textbf{laterolaterale}.

\subsection{Superfici del corpo}

Nel \textbf{corpo} si riconoscono: una superficie anteriore (ventrale), una posteriore (dorsale) e due laterali (destra e sinistra).

Negli \textbf{arti} si riconoscono: una superficie anteriore (ventrale), una posteriore (dorsale), una laterale e una mediale.

\subsection{Punti e linee di repere}

I \textbf{punti e linee di repere} (dal francese \textit{repère} = trovare) sono strutture scheletriche apprezzabili dall'esterno durante l'esplorazione superficiale del corpo, che permettono di localizzare la posizione di punti e aree rispetto a riferimenti fissi. Principali punti di repere cervicali e toracici superiori:

\begin{itemize}
  \item Processo trasverso dell'atlante
  \item Linea nucale posteriore superiore
  \item Processo trasverso della II vertebra cervicale
  \item Processo spinoso della VII vertebra cervicale
  \item I costa
\end{itemize}

% ---------------------------------------------------------------------------
\section{Le cavità del corpo}

Le \textbf{cavità del corpo} contengono e proteggono gli organi interni e permettono i cambiamenti di dimensione e forma dei visceri. Si suddividono in:

\begin{itemize}
  \item \textbf{Cavità dorsale}:
  \begin{itemize}
    \item \textit{Cavità cranica}: contiene l'encefalo.
    \item \textit{Cavità spinale}: contiene il midollo spinale.
  \end{itemize}
  \item \textbf{Cavità ventrale}:
  \begin{itemize}
    \item \textit{Cavità toracica}: contiene il mediastino (trachea, bronchi, esofago, timo, cavità pericardica con il cuore) e le cavità pleuriche (con i polmoni). Il \textbf{diaframma} separa la cavità toracica da quella addominopelvica.
    \item \textit{Cavità addominopelvica}:
    \begin{itemize}
      \item \textbf{Cavità addominale}: fegato, cistifellea, stomaco, intestino tenue e crasso, milza, reni e ureteri.
      \item \textbf{Cavità pelvica}: vescica urinaria, organi dell'apparato genitale e parte dell'intestino crasso.
    \end{itemize}
  \end{itemize}
\end{itemize}

La cavità addominopelvica viene convenzionalmente suddivisa in \textbf{4 quadranti} (superiore destro e sinistro, inferiore destro e sinistro) oppure in \textbf{9 regioni}: epigastrica, ipocondriaca destra e sinistra, ombelicale, lombare destra e sinistra, ipogastrica, iliaca destra e sinistra.

% ---------------------------------------------------------------------------
\section{L'apparato muscolare}

L'\textbf{apparato locomotore} integra tre sistemi: il \textbf{sistema scheletrico} (ossa che formano lo scheletro), il \textbf{sistema articolare} (articolazioni che mettono in giunzione le ossa) e il \textbf{sistema muscolare} (muscolatura scheletrica).

Il \textbf{muscolo} è un tessuto composto da fibre muscolari dotate di capacità contrattile. Le funzioni dei muscoli comprendono: movimento dello scheletro, mantenimento della postura, sostegno a parti molli, funzione protettiva e regolazione degli orifizi.

\subsection{Capacità del muscolo}

\begin{itemize}
  \item \textbf{Contrattilità}: capacità di accorciarsi o allungarsi.
  \item \textbf{Eccitabilità}: capacità di rispondere a stimoli nervosi e/o ormonali, permettendo al sistema nervoso e al sistema endocrino di regolare l'attività muscolare.
  \item \textbf{Estensibilità}: capacità di essere stirati fino alla loro lunghezza a riposo dopo la contrazione.
  \item \textbf{Elasticità}: capacità di ritornare alla lunghezza iniziale dopo una fase di stiramento.
\end{itemize}

\subsection{Tessuto muscolare}

Esistono due grandi tipologie di muscoli:

\begin{itemize}
  \item \textbf{Muscolatura striata} (``rossa'' o volontaria): regolata dalla volontà del soggetto. Le cellule sono lunghe, cilindriche, striate e multinucleate; la contrazione è attivata da motoneuroni.
  \item \textbf{Muscolatura liscia} (``bianca'' o involontaria): attività contrattile autonoma e indipendente; produce una contrazione prolungata.
\end{itemize}

Il \textbf{muscolo cardiaco} (miocardio) ha caratteristiche istologiche simili a quelle dei muscoli striati volontari, ma funziona come un muscolo involontario con contrazione ritmica spontanea.

\subsection{Classificazione dei muscoli}

\paragraph{In base al numero di punti di origine (capi):}

\begin{itemize}
  \item \textbf{Monocipiti}: un solo punto di origine.
  \item \textbf{Bicipiti}: 2 punti di origine.
  \item \textbf{Tricipiti}: 3 punti di origine.
  \item \textbf{Quadricipiti}: 4 punti di origine.
\end{itemize}

I \textbf{muscoli scheletrici} hanno sia l'origine che l'inserzione nelle ossa. I \textbf{muscoli pellicciai} hanno almeno un punto di attacco nel derma; la loro contrazione muove la cute.

\paragraph{In base alla forma:}

\begin{itemize}
  \item \textbf{M.\ lunghi}: la lunghezza prevale sulla larghezza e sullo spessore.
  \item \textbf{M.\ larghi}: lo spessore è nettamente inferiore alla lunghezza e alla larghezza.
  \item \textbf{M.\ brevi}: lunghezza, larghezza e spessore sono pressoché uguali.
  \item \textbf{M.\ anulari}: circondano gli orifizi naturali del corpo.
  \item \textbf{M.\ orbicolari}: si comportano come gli altri muscoli scheletrici.
  \item \textbf{Sfinteri}: hanno un accentuato tono muscolare e sono in continua contrazione.
\end{itemize}

\paragraph{In base all'orientamento delle fibre:}

\begin{itemize}
  \item \textbf{M.\ unipennati}: un tendine centrale su cui le fibre si inseriscono da un lato (es.\ vasto mediale, vasto laterale).
  \item \textbf{M.\ semipennati}: le fibre si inseriscono su due linee di attacco lineari e contrapposte (es.\ flessore lungo del pollice).
  \item \textbf{M.\ bipennati}: le fibre convergono da due diverse linee di origine sulle due facce di un tendine centrale (es.\ retto del femore, gastrocnemio).
  \item \textbf{M.\ multi- o pluripennati}: vari fasci tendinei con origine comune, sui quali si inseriscono diversi gruppi di fibre (es.\ deltoide).
\end{itemize}

\paragraph{In base al numero di punti di inserzione:}

\begin{itemize}
  \item \textbf{Monocaudati}: un solo punto di inserzione.
  \item \textbf{Bicaudati}: 2 punti di inserzione.
  \item \textbf{Tricaudati}: 3 punti di inserzione.
  \item \textbf{Pluricaudati}: più punti di inserzione.
\end{itemize}

\paragraph{In base alle iscrizioni tendinee:}

\begin{itemize}
  \item \textbf{Muscoli monogastrici}: nessun tendine intermedio (es.\ sopraioidei).
  \item \textbf{Muscoli digastrici}: un tendine intermedio.
  \item \textbf{Muscoli poligastrici}: più tendini intermedi (es.\ retto dell'addome).
\end{itemize}

\paragraph{In base alla funzione:}

Flessori ed estensori, adduttori e abduttori, pronatori e supinatori, rotatori interni ed esterni. In base al ruolo durante il movimento: \textbf{muscoli motori primari}, \textbf{muscoli sinergici}, \textbf{muscoli antagonisti}, \textbf{muscoli stabilizzatori} e \textbf{muscoli neutralizzatori}.

\subsection{L'angolo di pennazione}

L'\textbf{angolo di pennazione} è l'inclinazione delle fibre muscolari rispetto all'asse di trazione del tendine.

\begin{itemize}
  \item I \textbf{muscoli a fasci paralleli} (fusiformi, es.\ bicipite) trasmettono tutta la loro capacità contrattile al tendine.
  \item I \textbf{muscoli a fibre pennate} ne trasmettono solo una parte: un angolo di pennazione di 30° trasmette al tendine circa l'87\% della tensione (\(\cos 30^\circ = 0{,}866\)).
\end{itemize}

I muscoli pennati hanno più fibre compatte rispetto ai muscoli fusiformi, e la \textbf{forza del muscolo aumenta all'aumentare dell'angolo di pennazione}.

In sintesi, tre fattori determinano le capacità di forza di un muscolo:
\begin{itemize}
  \item La \textbf{forma}: determina il grado di contrazione e la quantità di forza sviluppabile.
  \item L'\textbf{ipertrofia}: 1 cm\textsuperscript{2} di sezione trasversa genera circa 2--3 kg di forza.
  \item L'\textbf{orientamento} delle fibre: ruolo chiave attraverso le leve articolari.
\end{itemize}

\noindent\textbf{Implicazioni per i calciatori:} un volume eccessivo dei quadricipiti e dei muscoli femorali, fatto crescere in modo rapido senza un accurato adattamento funzionale, può essere causa di gravi infortuni. Quando un atleta raggiunge un limite del volume muscolare, i muscoli pennati subiscono un aumento dell'angolo di pennazione che limita le possibilità di espressione di forza. Analogamente, l'allenamento con i pesi modifica l'architettura muscolare aumentando l'angolo di pennazione.

\subsection{La contrazione muscolare}

\begin{quote}
\textit{«Il muscolo esercita una forza dipendente dal numero delle sue fibre che è in rapporto con il lavoro che il muscolo può svolgere. Il lavoro muscolare è la risultante del prodotto del peso che un muscolo può sollevare per lo spostamento che può eseguire.»}
\end{quote}

I tre tipi fondamentali di contrazione muscolare sono:

\begin{itemize}
  \item \textbf{Concentrica}: il muscolo si accorcia mentre genera forza (es.\ fase di sollevamento del peso).
  \item \textbf{Eccentrica}: il muscolo si allunga controllando la forza (es.\ fase di abbassamento del peso).
  \item \textbf{Isometrica}: il muscolo si contrae ma non si accorcia, senza spostamento articolare.
\end{itemize}

\subsection{Le leve articolari}

Una \textbf{leva} è una struttura rigida che si muove facendo perno su un punto fisso (\textbf{fulcro}). Nelle leve anatomiche il fulcro corrisponde all'\textbf{articolazione}, il \textbf{braccio della forza} è la parte della leva tra il fulcro e il punto di applicazione della forza muscolare, e il \textbf{braccio del carico} è la parte tra il fulcro e il punto di applicazione della resistenza.

Si distinguono tre classi di leve:

\begin{itemize}
  \item \textbf{Leva di primo tipo}: il fulcro è interposto tra la forza applicata e la resistenza (es.\ equilibrio della testa sul rachide cervicale: i muscoli posteriori del collo sono la forza applicata, il peso del cranio è la resistenza, l'articolazione atlanto-occipitale è il fulcro).
  \item \textbf{Leva di secondo tipo}: fulcro e forza motrice si trovano dalla stessa parte, con la forza motrice più distante dal fulcro rispetto alla resistenza. È sempre \textbf{vantaggiosa} (es.\ alzarsi in punta di piedi: il fulcro è l'articolazione metatarso-falangea, la resistenza è il peso corporeo, la forza è il tricipite surale).
  \item \textbf{Leva di terzo tipo}: la forza si applica in un punto più vicino al fulcro rispetto alla resistenza. È la configurazione più comune nell'apparato locomotore umano (es.\ flessione dell'avambraccio: il fulcro è il gomito, la forza è l'inserzione del bicipite, la resistenza è il peso dell'avambraccio e dell'oggetto tenuto in mano).
\end{itemize}

\subsection{La cellula muscolare}

La \textbf{cellula muscolare} (fibrocellula muscolare) presenta le seguenti caratteristiche:

\begin{itemize}
  \item Forma \textbf{fusiforme}.
  \item Si accorcia in seguito a stimolo.
  \item Il citoplasma prende il nome di \textbf{sarcoplasma}.
  \item La membrana cellulare è detta \textbf{sarcolemma}.
  \item Contiene \textbf{miofibrille} composte dalle proteine contrattili \textbf{actina} e \textbf{miosina}.
\end{itemize}
